\clearpage\fancyhead[L]{\small GAUSS AN VON LINDENAU.}\SourceLeaf{321}

\begin{center}\textbf{GAUSS AN VON LINDENAU.}\par\vspace{4pt}\hrule\vspace{4pt}\fontsize{10}{12}\selectfont Monatliche Correspondenz. Band. XXI. S. 276\,..\,280. 1810 März.\par\vspace{4pt}\hrule\vspace{4pt}\end{center}
\begin{flushright}Göttingen, den 23. Febr. 1810.\end{flushright}

\par\indent ---\enspace---\enspace--- Ich habe so eben einige Rechnungen über die beiden letzten Oppositionen der Pallas beendigt, deren Resultate ich in Nr. 32 der Göttinger gelehrten Zeitungen [S. 317 d. B.] habe abdrucken lassen und die ich Ihnen hier mittheile. ---\enspace---

Ich habe auch schon einige andere Resultate in Rücksicht der Elemente darauf gegründet, und bin jetzt mit noch einigen Rechnungen darüber beschäftigt: dies zusammen, nebst einer Auseinandersetzung verschiedener Kunstgriffe, welche ich schon seit vielen Jahren bei Anwendung der im 3. Abschnitt des 2ten Buches meiner \emph{Theoria} entwickelten Methode (\emph{des moindres carrés}) gebraucht habe, bestimme ich zu einem Aufsatz für unsere Societät, aus welchem ich Ihnen, sobald die Arbeit vollendet ist, einen Auszug schicken werde.

Die hiesigen Beobachtungen der \emph{Vesta}, welche alle vom Prof. \textsc{Harding} am Mauerquadranten gemacht sind, füge ich hier bei.

\Needspace{12\baselineskip}\begin{center}\emph{Vesta Beobachtungen.}\end{center}
\begin{center}\fontsize{10}{12}\selectfont\setlength{\tabcolsep}{5pt}\renewcommand{\arraystretch}{1.15}\begin{tabular}{lrrr}
1810 & \shortstack{Mittlere Zeit\\in Göttingen} & \shortstack{Scheinbare AR.\\der Vesta} & \shortstack{Scheinbare\\nördl. Declin.}\\
Januar 13 & $11^{\text{u}}\;1^{\text{m}}\;27^{\text{s}}7$ & $98^\circ\;5'\;44''5$ & $23^\circ\;19'\;17''6$\\
15 & $10\;51\;30.1$ & $97\;34\;57.0$ & $23\;26\;14.3$\\
28 & $9\;49\;13.7$ & $94\;46\;15.2$ & $24\;7\;33.7$\\
Februar 2 & $9\;26\;32.7$ & $94\;0\;51.3$ & $24\;20\;51.7$\\
3 & $9\;22\;7.4$ & $93\;53\;19.2$ & $24\;23\;37.8$\\
6 & $9\;8\;56.1$ & $93\;32\;33.2$ & $24\;31\;1.9$\\
\end{tabular}\end{center}

\par\indent Hier sind auch unsere Beobachtungen der Bedeckung des Jupiter vom 8. Febr. Die Austritte sind ganz unsicher, weil der Mond sehr niedrig stand.

\Needspace{17\baselineskip}\begin{center}\emph{Eintritte.}\end{center}
\begin{center}\fontsize{10}{12}\selectfont\setlength{\tabcolsep}{5pt}\renewcommand{\arraystretch}{1.15}\begin{tabular}{lrl}
1. Trabant & $8^{\text{u}}\;58^{\text{m}}\;2^{\text{s}}9$ & \textsc{Harding}\\
 & $3.3$ & \textsc{Gauss}\\
2. Trabant & $8\;58\;26.2$ & G.\\
 & $28.7$ & H.\\
$\JupiterSymbol$ I. R. & $9\;0\;16.4$ & H.\\
 & $17.4$ & G.\\
$\JupiterSymbol$ II. R. & $9\;2\;26.0$ & G.\\
 & $26.7$ & H.\\
3. Trabant & $9\;7\;55.5$ & G.\\
 & $56.4$ & H.\\
4. Trabant & $9\;14\;35.3$ & H.\\
 & $37.0$ & G.\\
\end{tabular}\end{center}
\Needspace{9\baselineskip}\begin{center}\emph{Austritte.}\end{center}
\begin{center}\fontsize{10}{12}\selectfont\setlength{\tabcolsep}{5pt}\renewcommand{\arraystretch}{1.15}\begin{tabular}{lrl}
1. Trabant & $9^{\text{u}}\;30^{\text{m}}\;59^{\text{s}}$ & \textsc{Gauss} schon ausgetreten\\
$\JupiterSymbol$ I. R. & $9\;33\;37.9$ & H.\\
II. R. & $9\;35\;19.6$ & G.\\
 & $36.1$ & H.\\
3. Trabant & $9\;43\;39.3$ & H.\\
\end{tabular}\end{center}

\par\indent Für die Notirung der Druckfehler in meiner Theoria, bin ich Herrn \textsc{Oriani} sehr verbunden. Er hat ganz Recht, dass ich pag. 129 hinzu zu fügen vergessen habe, dass $B, B', B'' = 0$, vorausgesetzt
\SourceLeaf{322}werden müssen, wenn die Bedingungsgleichung, bei welcher die Gleichung [7] unbrauchbar ist, die dort angegebene Gestalt haben soll. Es ist übrigens klar, dass wenn auch nicht $B, B', B''=0$ sind, doch die Gleichung [7] unbrauchbar sein kann, wenn nemlich  der 12 gliedrige Ausdruck, welchen \textsc{Oriani} entwickelt hat, zufällig $=0$ oder sehr klein wird. Dass \textsc{Euler} schon das Theorem gefunden hat, woraus der schöne von mir \textsc{La Place} beigelegte Lehrsatz sehr leicht abgeleitet werden kann, fiel mir selbst schon früher ein, als aber die Stelle pag. 212 schon abgedruckt war: ich wollte es aber nicht unter die Errata setzen, weil \textsc{La Place} wenigstens das obige Theorem doch erst in der dort gebrauchten Form aufgestellt hat. Die meisten der von \textsc{Oriani} angezeigten Druckfehler hatte ich mir auch schon notirt. Hier sind noch drei andere von ihm übersehene:

\begin{center}\fontsize{10}{12}\selectfont\setlength{\tabcolsep}{5pt}\renewcommand{\arraystretch}{1.15}\begin{tabular}{lrlrlllll}
Pag. & 1 & Zeile & 4 & v. u. & statt & inversa & l. & composita\\
--- & 65 & --- & 2 & --- & --- & $\cos\theta$. & l. & $\cos\zeta$\\
--- & 195 & --- & 16 &  & --- & $\lambda'''$ & l. & $=\lambda'''$.\\
\end{tabular}\end{center}

\par\indent Den Druckfehler $380^\circ$ statt $180^\circ$ S. 4 finde ich in meinem Exemplare nicht. Statt der zwei Druckfehler pag. 83 könnte einer nemlich $\zeta+45$ statt $\zeta-45^\circ$ gesetzt werden.
\clearpage\fancyhead[L]{\small PALLAS.}

\begin{center}\textbf{PALLAS.}\par\vspace{4pt}\hrule\vspace{4pt}\fontsize{10}{12}\selectfont Göttingische gelehrte Anzeigen. Stück 32. Seite 305\,..\,308. 1810 Februar 24.\par\vspace{4pt}\hrule\vspace{4pt}\end{center}

\par\indent Die \emph{Pallas} war im Jahre 1808 nur auf wenigen Sternwarten, und nur unvollkommen, beobachtet: daher auch seit 1807 keine Berichtigung der Elemente hatte unternommen werden können. Auf der hiesigen Sternwarte hatte die grosse Lichtschwäche des Planeten alle Beobachtungen am Mauerquadranten unmöglich gemacht, und die wenigen am Kreismicrometer angestellten Beobachtungen konnten, besonderer Umstände wegen, auf grosse Genauigkeit keinen Anspruch machen. Auf der Seeberger Sternwarte beobachtete zwar Hr. \textsc{von Lindenau} den Planeten am Passageninstrumente mit gewohnter Schärfe: allein Declinationen konnten auch dort nicht beobachtet werden. Die in Mailand am Aequatorealsector gemachten Bestimmungen fangen erst am 22. August an, und konnten daher nicht wohl mehr für die Opposition gebraucht werden. Erst durch das Astron. Jahrbuch für 1812 sind nun noch einige um die Zeit der Opposition zu Prag von Hrn. \textsc{David} angestellte vollständige Beobachtungen bekannt geworden, von denen Hr. Prof. \textsc{Gauss} die Declinationen benutzt hat, um daraus, in Verbindung mit den auf der Seeberger Sternwarte bestimmten geraden Aufsteigungen die Opposition abzuleiten. Leider hat die Untersuchung jener Prager Beobachtungen gezeigt, dass sie schlecht harmoniren, da z.\,B. die verschiedenen Declinationen Resultate geben, welche beinahe Eine Minute von einander abweichen. Da es inzwischen an bessern Declinationen durchaus fehlt, so blieb nichts übrig, als das, was da war, so gut wie möglich, zu benutzen, und so ergab sich für die Opposition von 1808, die fünfte, welche bisher beobachtet wurde, folgendes Resultat:

\begin{center}Opposition der Pallas 1808 Juli 26.\\$21^{\text{u}}\;17^{\text{m}}\;32^{\text{s}}$ Meridian von Göttingen\par Wahre Länge .\;.\;.\;.\;.\; $304^\circ\;2'\;59''7$\par Wahre geocentrische Breite\quad $37\;33\;54$ Nordlich\end{center}

\SourceLeaf{323}\par\indent Im Jahre 1809 scheinen die auf der hiesigen Sternwarte gemachten Beobachtungen der Pallas (s. Gel. Anz. 1809 S. 1348) [1809 Aug. 26 Seite 314 dieses Bandes] die frühesten gewesen zu sein. Von auswärtigen Sternwarten hat Hr. Prof. \textsc{Gauss} durch die gefällige Mittheilung der Herren \textsc{von Lindenau} und \textsc{Bouvard} eine schätzbare Reihe Seeberger und Pariser Meridianbeobachtungen erhalten. Die letztern, die sonst erst nach einigen Jahren bekannt werden, lassen wir hier folgen: die mittlern Zeiten, welche von Hrn. \textsc{Bouvard} nicht beigefügt waren, hat Hr. Prof. \textsc{Gauss} selbst hinzugesetzt.

\Needspace{18\baselineskip}\begin{center}\emph{Beobachtungen der Pallas zu Paris.}\end{center}
\begin{center}\fontsize{9}{11}\selectfont\setlength{\tabcolsep}{5pt}\renewcommand{\arraystretch}{1.15}\begin{tabular}{lrrr}
1809 & Mittlere Zeit & \shortstack{Scheinbare gerade\\Aufsteigung} & \shortstack{Scheinbare südl.\\Abweichung}\\
Sept. 12 & $12^{\text{u}}\;51^{\text{m}}\;39^{\text{s}}6$ & $4^\circ\;30'\;25''80$ & $4^\circ\;20'\;10''4$\\
15 & $12\;37\;41.7$ & $3\;57\;48.75$ & $5\;5\;16.4$\\
28 & $11\;36\;52.6$ & $1\;31\;50.25$ & $8\;20\;23.4$\\
29 & $11\;32\;11.1$ & $1\;20\;24.30$ & $8\;34\;57.3$\\
Octob. 2 & $11\;18\;7.1$ & $0\;46\;15.30$ & $9\;18\;22.5$\\
5 & $11\;4\;5.9$ & $0\;12\;35.25$ & $10\;0\;21.0$\\
6 & $10\;59\;25.5$ & $0\;1\;37.50$ & $10\;14\;6.9$\\
7 & $10\;54\;46.2$ & $359\;50\;44.55$ & $10\;27\;36.8$\\
8 & $10\;50\;7.1$ & $359\;39\;54.90$ & $10\;40\;1.5$\\
9 & $10\;45\;28.7$ & $359\;29\;16.50$ & $10\;54\;5.5$\\
13 & $10\;27\;2.2$ & $358\;48\;26.25$ & $11\;54\;59.5$\\
\end{tabular}\end{center}

\par\indent Die hier abgedruckten Zahlen stimmen genau mit den Angaben in Hrn. \textsc{Bouvard}'s Schreiben überein: allein die Discussion dieser Bestimmung hat gezeigt, dass die gerade Aufsteigung vom 12. Sept. um etwa $1'\,15''$ in Bogen zu gross (die geraden Aufsteigungen waren blos in Zeit angegeben); die Declination vom 8. Oct. um 1 Minute zu klein, die Declination vom 13. Oct. um 10 Minuten zu gross ist. Das Resultat aus diesen Beobachtungen sind folgende Bestimmungen für die sechste Opposition der Pallas:

\begin{center}Zeit der Opposition 1809 Sept. 22.\\$16^{\text{u}}\;10^{\text{m}}\;20^{\text{s}}$ Meridian von Göttingen\par Wahre Länge der Pallas .\;.\; $359^\circ\;40'\;4''4$\par Wahre geocentrische Breite ,\quad $7\;22\;10.1$ südlich\end{center}

\par\indent Die von Hrn. Prof. \textsc{Gauss} aus diesen beiden Oppositionen, in Verbindung mit den vier frühern, abgeleiteten Resultaten werden bei einer andern Gelegenheit bekannt gemacht werden.

Wir lassen hier noch, eben so wie im vorigen Jahre, ein uns unlängst zur Bekanntmachung zugesandtes Verzeichniss einiger im Jahre 1810 sichtbarer und in den astronomischen Ephemeriden entweder gar nicht oder nicht richtig angezeigter Sternbedeckungen, welche Hr. \textsc{Wisniewsky} für den Meridian und Horizont von Berlin im voraus berechnet hat, abdrucken. --- ---
